\documentclass[12pt]{article}
\usepackage{amsmath, amssymb, amsthm}
\usepackage{geometry}
\usepackage{enumitem}
\usepackage{hyperref}
\usepackage{bm}
\usepackage{mathtools}

\geometry{margin=1in}

\title{\textbf{ORF 309 --- Slides Relevant to Homework 3, Question 4}\\
\large Compiled from Slides 05, 06, 07}
\author{}
\date{}

\newtheoremstyle{definition_style}{6pt}{6pt}{}{}{\bfseries}{.}{.5em}{}
\theoremstyle{definition_style}
\newtheorem{defn}{Definition}[section]
\newtheorem{thm}[defn]{Theorem}
\newtheorem{claim}[defn]{Claim}
\newtheorem{example}[defn]{Example}
\newtheorem{remark}[defn]{Remark}
\newtheorem{exercise}[defn]{Exercise}

\begin{document}
\maketitle
\tableofcontents
\newpage

%=============================================================
\section{Homework 3, Question 4 --- The Problem}
%=============================================================

\textbf{Q.4 (Love triangles).} A dating website works as follows. There are $n$ people signed up.
For every pair of people, the website independently flips a coin that comes up heads with probability $p$.
If the coin comes up heads, the website creates a match between this pair of people.
(Being a scam, the site does not actually use any specific information about the people in deciding who to hook up.)

In this setting, it is perfectly possible that the website randomly creates a \emph{love triangle}: three people who are all mutually hooked up.

\begin{enumerate}[label=\textbf{\alph*)}]
    \item What is the expected number of hookups a given individual has at once?
    \item What is the probability that a given individual has more than one hookup?
    \item What is the expected number of love triangles on the entire site?
\end{enumerate}

\medskip
\noindent\textbf{Key ideas needed to solve each part:}
\begin{itemize}
    \item \textbf{Part (a):} Indicator random variables + linearity of expectation. Each potential partner gives an independent Bernoulli trial; sum the indicators.
    \item \textbf{Part (b):} The number of hookups for a given person follows $\mathrm{Binom}(n-1, p)$; use the Binomial PMF to compute $P\{\text{more than 1 hookup}\} = 1 - P\{0\} - P\{1\}$.
    \item \textbf{Part (c):} Use indicator random variables over all $\binom{n}{3}$ triples; apply linearity of expectation and independence of edges.
\end{itemize}

\newpage
%=============================================================
\section{Discrete Random Variables and Expectations (Slides 05)}
%=============================================================

\subsection{Types of Random Variables}

\begin{defn}[Discrete Random Variable]
A random variable $X: \Omega \to D$ is called \textbf{discrete} if the value set $D$ is a finite or countable set.
\end{defn}

\begin{defn}[Continuous Random Variable]
A random variable $X: \Omega \to D$ is called \textbf{continuous} if the value set $D$ is NOT countable.
\end{defn}

\subsection{Expectation}

\begin{defn}[Expectation of a Discrete Random Variable]
If $X: \Omega \to D$ is a discrete r.v.\ (DRV) and $f: D \to \mathbb{R}$ is a value function, then the \textbf{expectation} of $f(X)$ is defined as
\[
    E[f(X)] = \sum_{i \in D} f(i) \, P\{X = i\}.
\]
If $D \subset \mathbb{R}$, we can choose $f(i) = i$ and the expectation of $X$ becomes:
\[
    E[X] = \sum_{i \in D} i \, P\{X = i\}.
\]
\end{defn}

\begin{remark}
This definition makes sense only for DRVs, because for a continuous r.v.\ (CRV) we have $P\{X = x\} = 0$ for all $x \in D$.
\end{remark}

\subsection{The Indicator Function (Very Important!)}

\begin{defn}[Indicator Random Variable]
Assume $A$ is an event. Then we can define the DRV $\mathbf{1}_A$ with values:
\[
    \mathbf{1}_A(\omega) = \begin{cases} 1, & \text{if } \omega \in A \\ 0, & \text{if } \omega \notin A \end{cases}
\]
The DRV $\mathbf{1}_A$ is called the \textbf{indicator random variable} for event $A$.
\end{defn}

\begin{thm}[Expectation of the Indicator]
\[
    E[\mathbf{1}_A] = 1 \cdot P\{\mathbf{1}_A = 1\} + 0 \cdot P\{\mathbf{1}_A = 0\} = P\{A\}.
\]
That is, \textbf{the expectation of $\mathbf{1}_A$ equals the probability of event $A$}.
\end{thm}

\subsection{Usefulness of the Indicator Function}

The indicator function allows us to represent the number of successes in a series of trials as a simple sum.

\begin{example}[Coin Flips]
Consider the DRV $X$ representing the number of heads in $n$ coin flips.
\begin{itemize}
    \item Define the events: $A_i = \{\text{the } i\text{-th flip is heads}\}$ for $i = 1, 2, \ldots, n$.
    \item For each event $A_i$, define the indicator function $\mathbf{1}_{A_i}$.
    \item Then we can express $X$ as the sum of indicators:
    \[
        X = \mathbf{1}_{A_1} + \mathbf{1}_{A_2} + \cdots + \mathbf{1}_{A_n}.
    \]
\end{itemize}
This representation is the key trick for computing expectations of counts.
\end{example}

\subsection{Distributions of DRVs}

\begin{defn}[Distribution]
If $X: \Omega \to D$ is a DRV, then its \textbf{distribution} is given by specifying the probabilities of all its outcomes $\{P\{X = i\}\}_{i \in D}$.
\end{defn}

\begin{defn}[Uniform Distribution]
Let $D$ be a finite set. The DRV $X: \Omega \to D$ for which
\[
    P\{X = i\} = \frac{1}{|D|}, \quad \forall i \in D
\]
is said to have the \textbf{Uniform Distribution} on $D$. Notation: $X \sim \mathrm{Unif}(D)$.
\end{defn}

\subsection{Joint and Marginal Distributions}

\begin{defn}[Joint Distribution]
The \textbf{joint distribution} of two DRVs $X: \Omega \to D_X$ and $Y: \Omega \to D_Y$ is specified by the probabilities
\[
    \{P\{X = i, Y = j\}\}_{i \in D_X,\, j \in D_Y}.
\]
\end{defn}

\begin{defn}[Marginal Distribution]
The distribution of DRV $X$ only is called the \textbf{marginal distribution} of $X$.
\end{defn}

\begin{claim}[Recovering Marginals from Joint]
\[
    P\{X = i\} = \sum_{j \in D_Y} P\{X = i, Y = j\}.
\]
\textit{Proof sketch:} Write $\Omega = \bigcup_{j \in D_Y}\{Y = j\}$ (mutually exclusive partition), so $\{X=i\} = \bigcup_{j \in D_Y}\{X=i, Y=j\}$, and take probabilities.
\end{claim}

\subsection{Linearity of Expectation}

\begin{thm}[Linearity of Expectation]
For any two DRVs $X$ and $Y$ (no independence required):
\[
    E[X + Y] = E[X] + E[Y].
\]
\end{thm}

\begin{proof}
\begin{align*}
E[X + Y] &= \sum_{x \in D_X} \sum_{y \in D_Y} (x + y) \, P\{X = x, Y = y\} \\
          &= \sum_{x \in D_X} \sum_{y \in D_Y} x \, P\{X=x, Y=y\} + \sum_{x \in D_X} \sum_{y \in D_Y} y \, P\{X=x, Y=y\} \\
          &= \sum_{x \in D_X} x \underbrace{\sum_{y \in D_Y} P\{X=x, Y=y\}}_{= P\{X=x\}} + \sum_{y \in D_Y} y \underbrace{\sum_{x \in D_X} P\{X=x, Y=y\}}_{= P\{Y=y\}} \\
          &= E[X] + E[Y].
\end{align*}
\end{proof}

\begin{remark}
By induction, linearity extends to any finite number of random variables:
\[
    E[X_1 + X_2 + \cdots + X_n] = E[X_1] + E[X_2] + \cdots + E[X_n].
\]
\textbf{This holds regardless of whether the $X_i$ are independent.}
\end{remark}

\subsection{Indicator Function Properties (Useful Identities)}

\begin{exercise}
Prove the following properties of indicator functions:
\begin{align}
    \mathbf{1}_{A \cap B} &= \mathbf{1}_A \cdot \mathbf{1}_B \\
    \mathbf{1}_{A \cup B} &= \mathbf{1}_A + \mathbf{1}_B - \mathbf{1}_{A \cap B} \\
    \mathbf{1}_{A^c}      &= 1 - \mathbf{1}_A
\end{align}
\end{exercise}

\newpage
%=============================================================
\section{Independence, Indicator Functions, and Conditional Expectation (Slides 06)}
%=============================================================

\subsection{Indicator Function and Expectation: The Key Example}

\begin{example}[{$n$ Coin Flips --- Proof that $E[X] = np$}]
\textbf{Setup:} Flip $n$ identical coins. Let $X = $ total number of heads.
\begin{itemize}
    \item Define events $A_i = \{\text{the } i\text{-th coin shows H}\}$, with $P\{A_i\} = p$ and $P\{A_i^c\} = 1-p$.
    \item Define indicator $\mathbf{1}_{A_i}$ for each $A_i$.
    \item Then $X = \mathbf{1}_{A_1} + \mathbf{1}_{A_2} + \cdots + \mathbf{1}_{A_n}$.
\end{itemize}
By linearity of expectation and $E[\mathbf{1}_{A_i}] = P\{A_i\} = p$:
\[
    E[X] = E[\mathbf{1}_{A_1}] + E[\mathbf{1}_{A_2}] + \cdots + E[\mathbf{1}_{A_n}] = p + p + \cdots + p = np.
\]

\textbf{Key remark:} We did \emph{not} need to assume $A_i$ and $A_j$ are independent for any $i \ne j$. Even if the coins were correlated (e.g., through collusion), the expected total heads is still $np$. Independence is not required for linearity of expectation.
\end{example}

\subsection{Independence of DRVs}

\begin{defn}[Independence of Two DRVs]
Two DRVs $X: \Omega \to D_X$ and $Y: \Omega \to D_Y$ are called \textbf{independent} ($X \perp Y$) if for every $x \in D_X$ and $y \in D_Y$:
\[
    P\{X = x \mid Y = y\} = P\{X = x\} \quad \Longleftrightarrow \quad P\{X = x, Y = y\} = P\{X = x\} \cdot P\{Y = y\}.
\]
\end{defn}

\begin{remark}
If $X \perp Y$, then their marginal distributions \emph{determine} their joint distribution.
\end{remark}

\begin{thm}[Key Rules]
For any DRVs $X$ and $Y$:
\begin{itemize}
    \item $E[X + Y] = E[X] + E[Y]$ \quad (\textbf{always}, no independence required).
    \item $E[\varphi_1(X) \cdot \varphi_2(Y)] = E[\varphi_1(X)] \cdot E[\varphi_2(Y)]$ \quad (\textbf{only if} $X \perp Y$).
\end{itemize}
\end{thm}

\begin{defn}[Independence of $n$ DRVs]
The $n$ DRVs $X_1, X_2, \ldots, X_n$ are called \textbf{independent} iff for all subsets $\{j_1, j_2, \ldots, j_k\} \subseteq \{1,\ldots,n\}$ and all values:
\[
    P\{X_{j_1} = x_{j_1}, X_{j_2} = x_{j_2}, \ldots, X_{j_k} = x_{j_k}\} = P\{X_{j_1} = x_{j_1}\} \cdot P\{X_{j_2} = x_{j_2}\} \cdots P\{X_{j_k} = x_{j_k}\}.
\]
\end{defn}

\subsection{Conditional Expectation}

\begin{defn}[Conditional Expectation]
If $X: \Omega \to D_X$ and $Y: \Omega \to D_Y$ are DRVs and $\varphi: D_X \to \mathbb{R}$, then the \textbf{conditional expectation} of $\varphi(X)$ given $Y = y$ is
\[
    E[\varphi(X) \mid Y = y] = \sum_{x \in D_X} \varphi(x) \, P\{X = x \mid Y = y\}.
\]
\end{defn}

\begin{example}[Die Rolling Until a 6]
\textbf{Setup:}
\begin{itemize}
    \item Roll a fair die repeatedly; stop when we see a 6.
    \item $T = $ the trial number when the first 6 appears.
    \item $X = $ the number of 1s seen before the game ends.
    \item Find: $E[X \mid T = k]$.
\end{itemize}
Define $Z_i = $ outcome of the $i$-th trial, with indicator $\mathbf{1}_{\{Z_i = 1\}}$.

If the game ends at trial $k$, then $X = \sum_{i=1}^{k-1} \mathbf{1}_{\{Z_i = 1\}}$, so:
\[
    E[X \mid T = k] = \sum_{i=1}^{k-1} P\{Z_i = 1 \mid T = k\}.
\]
Using independence of trials and $\{T = k\} = \{Z_1 \ne 6, \ldots, Z_{k-1} \ne 6, Z_k = 6\}$:
\[
    P\{Z_i = 1 \mid T = k\} = \frac{P\{Z_i = 1, Z_i \ne 6\}}{P\{Z_i \ne 6\}} = \frac{1/6}{5/6} = \frac{1}{5}, \quad \text{for all } i = 1, \ldots, k-1.
\]
Therefore:
\[
    E[X \mid T = k] = \sum_{i=1}^{k-1} \frac{1}{5} = \frac{k-1}{5}.
\]
\end{example}

\newpage
%=============================================================
\section{Bernoulli Processes and the Binomial Distribution (Slides 07)}
%=============================================================

\subsection{Bernoulli Random Variables and Processes}

\begin{defn}[Bernoulli Random Variable]
Any random variable $X: \Omega \to \{0, 1\}$ is called a \textbf{Bernoulli RV}.
\begin{itemize}
    \item $X = 1$ is a ``success''; $X = 0$ is a ``failure''.
    \item The indicator function $\mathbf{1}_A$ is a Bernoulli RV for every event $A$.
\end{itemize}
\end{defn}

\begin{defn}[Bernoulli Process]
A \textbf{Bernoulli Process} is an infinite sequence $X_1, X_2, X_3, \ldots$ of i.i.d.\ Bernoulli RVs with
\[
    P\{X_i = 1\} = p \quad \text{and} \quad P\{X_i = 0\} = 1 - p =: q.
\]
\end{defn}

\subsection{Counting Successes}

\begin{defn}[Counting RV $S_n$]
Define
\[
    S_n = X_1 + X_2 + \cdots + X_n,
\]
the \textbf{number of successes} in $n$ trials. Each $X_i \in \{0,1\}$ with $P\{X_i = 1\} = p$, $P\{X_i = 0\} = q$, and $p + q = 1$.
\end{defn}

\subsection{Deriving the Binomial Distribution}

\begin{example}[Exactly 2 Successes in 8 Trials]
We want $P\{S_8 = 2\}$.
\begin{itemize}
    \item Each specific sequence of 2 successes and 6 failures (e.g., $ssfffffff$) has probability $p^2 q^6$ by independence.
    \item The number of such sequences is $\binom{8}{2} = 28$ (choose which 2 of the 8 slots are successes).
    \item Therefore: $P\{S_8 = 2\} = \binom{8}{2} p^2 q^6$.
\end{itemize}
\end{example}

\textbf{Generalization:} For exactly $k$ successes in $n$ trials:
\[
    P\{S_n = k\} = \binom{n}{k} p^k q^{n-k}, \quad k = 0, 1, \ldots, n.
\]
where the binomial coefficient is
\[
    \binom{n}{k} = \frac{n!}{k!\,(n-k)!} = \frac{n(n-1)\cdots(n-k+1)}{k!}.
\]

\subsection{Binomial Distribution}

\begin{defn}[Binomial Distribution]
A random variable $Z$ with values $\{0, 1, 2, \ldots, n\}$ has the \textbf{Binomial Distribution} if
\[
    P\{Z = k\} = \binom{n}{k} p^k q^{n-k}, \quad k = 0, 1, \ldots, n,
\]
where $p \in (0,1)$ and $q = 1-p$. We write $Z \sim \mathrm{Binom}(n, p)$.

In particular, $S_n \sim \mathrm{Binom}(n, p)$.
\end{defn}

\begin{thm}[Sanity Check: Probabilities Sum to 1]
\[
    \sum_{k=0}^{n} P\{Z = k\} = \sum_{k=0}^{n} \binom{n}{k} p^k q^{n-k} = (p + q)^n = 1^n = 1. \quad \checkmark
\]
(Binomial theorem: $(p+q)^n = \sum_{k=0}^n \binom{n}{k} p^k q^{n-k}$.)
\end{thm}

\subsection{Expected Number of Successes}

\begin{thm}[{$E[S_n] = np$}]
If $S_n \sim \mathrm{Binom}(n, p)$, then $E[S_n] = np$.
\end{thm}

\begin{proof}[Proof via Indicator Functions (simpler)]
Write $S_n = X_1 + X_2 + \cdots + X_n$ where each $X_i$ is a Bernoulli$(p)$ r.v. By linearity of expectation:
\[
    E[S_n] = E[X_1] + E[X_2] + \cdots + E[X_n] = p + p + \cdots + p = np.
\]
\end{proof}

\begin{proof}[Proof via Algebra (alternative)]
\begin{align*}
    E[S_n] &= \sum_{k=1}^{n} k \binom{n}{k} p^k q^{n-k} \quad \text{(term $k=0$ vanishes)} \\
           &= \sum_{k=1}^{n} n \binom{n-1}{k-1} p^k q^{n-k} \quad \text{(since } k\binom{n}{k} = n\binom{n-1}{k-1}\text{)} \\
           &= np \sum_{\ell=0}^{n-1} \binom{n-1}{\ell} p^{\ell} q^{n-1-\ell} \quad \text{(substituting } \ell = k-1\text{)} \\
           &= np \cdot (p+q)^{n-1} = np.
\end{align*}
\end{proof}

\subsection{Successes Between Trials}

\begin{example}[Successes between trials 10 and 21]
Let $Z = S_{21} - S_9 = \sum_{k=10}^{21} X_k$. Since $X_{10}, X_{11}, \ldots, X_{21}$ are i.i.d.\ Bernoulli$(p)$ RVs, they have the same distribution as $X_1, \ldots, X_{12}$. Therefore:
\[
    S_{21} - S_9 \sim \mathrm{Binom}(12, p).
\]
\end{example}

\subsection{Incremental Successes (Independence of Non-Overlapping Intervals)}

\begin{example}[$P\{S_3 = 2, S_7 = 3\}$]
\begin{itemize}
    \item $S_3$ and $S_7$ are \textbf{not} independent (since $S_7$ depends on trials 1--3 which determine $S_3$).
    \item Rewrite: $\{S_3 = 2, S_7 = 3\} = \{S_3 = 2, S_7 - S_3 = 1\}$.
    \item $S_3 \sim \mathrm{Binom}(3, p)$ and $S_7 - S_3 \sim \mathrm{Binom}(4, p)$, and they \textbf{are} independent (non-overlapping intervals).
    \item Therefore:
    \[
        P\{S_3 = 2, S_7 = 3\} = P\{S_3 = 2\} \cdot P\{S_7 - S_3 = 1\} = \binom{3}{2} p^2 q \cdot \binom{4}{1} p q^3.
    \]
\end{itemize}
\end{example}

\newpage
%=============================================================
\section{Applying the Slides to HW3 Q4}
%=============================================================

This section shows how the tools above solve each part of Q4.

\subsection{Part (a): Expected Number of Hookups for a Given Person}

Fix a person $i$. There are $n-1$ other people; for each person $j \ne i$, independently the website flips a coin with $P\{\text{heads}\} = p$.

Define the indicator for person $i$ being matched with person $j$:
\[
    \mathbf{1}_{ij} = \begin{cases} 1, & \text{if persons } i \text{ and } j \text{ are matched} \\ 0, & \text{otherwise} \end{cases}
\]
with $E[\mathbf{1}_{ij}] = P\{\text{match}\} = p$.

The total hookups for person $i$ is:
\[
    H_i = \sum_{j \ne i} \mathbf{1}_{ij}.
\]
By linearity of expectation:
\[
    \boxed{E[H_i] = \sum_{j \ne i} E[\mathbf{1}_{ij}] = (n-1) \cdot p.}
\]

\subsection{Part (b): Probability of More than One Hookup}

The hookups for person $i$ are $H_i = \sum_{j \ne i} \mathbf{1}_{ij}$, where the $\mathbf{1}_{ij}$ are independent Bernoulli$(p)$ RVs (independent by assumption). Therefore:
\[
    H_i \sim \mathrm{Binom}(n-1, p).
\]
\begin{align*}
    P\{H_i > 1\} &= 1 - P\{H_i = 0\} - P\{H_i = 1\} \\
                 &= 1 - \binom{n-1}{0} p^0 q^{n-1} - \binom{n-1}{1} p^1 q^{n-2} \\
                 &= \boxed{1 - q^{n-1} - (n-1)p\,q^{n-2},}
\end{align*}
where $q = 1 - p$.

\subsection{Part (c): Expected Number of Love Triangles}

A love triangle is a set $\{i, j, k\}$ of three people who are all mutually matched.

\textbf{Step 1: Define an indicator for each triple.} For every unordered triple $\{i,j,k\}$, define:
\[
    \mathbf{1}_{ijk} = \begin{cases} 1, & \text{if all three pairs } (i,j),\, (j,k),\, (i,k) \text{ are matched} \\ 0, & \text{otherwise} \end{cases}
\]

\textbf{Step 2: Compute $E[\mathbf{1}_{ijk}]$.} The three matches are independent (separate coin flips), so:
\[
    E[\mathbf{1}_{ijk}] = P\{\text{all three pairs matched}\} = p \cdot p \cdot p = p^3.
\]

\textbf{Step 3: Count all possible triples.} The number of ways to choose 3 people from $n$ is $\binom{n}{3} = \frac{n(n-1)(n-2)}{6}$.

\textbf{Step 4: Apply linearity of expectation.}
\[
    E[\text{total love triangles}] = \sum_{\{i,j,k\}} E[\mathbf{1}_{ijk}] = \binom{n}{3} \cdot p^3 = \boxed{\frac{n(n-1)(n-2)}{6} \cdot p^3.}
\]

\end{document}
